% ============================================================
% 06-web-sucelje.tex — PRVA VERZIJA
% ============================================================

Razvijeno je i web su\v{c}elje za interaktivnu
vizualizaciju procesa generiranja sinteti\v{c}kih primjera razli\v{c}itim
SMOTE varijantama. Su\v{c}elje omogu\'{c}uje korisniku da na intuitivan
na\v{c}in istra\v{z}i razlike me\dj{}u algoritmima i utjecaj njihovih
parametara na distribuciju generiranih uzoraka.

\subsubsection{Motivacija}

Razumijevanje na\v{c}ina na koji razli\v{c}iti SMOTE algoritmi generiraju
sinteti\v{c}ke primjere \v{c}esto je ote\v{z}ano zbog vi\v{s}edimenzionalne
prirode podataka. Vizualizacija u dvije ili tri dimenzije, dobivena
redukcijom dimenzionalnosti, omogu\'{c}uje intuitivno sagledavanje
razlika izme\dj{}u linearnih pristupa (klasi\v{c}ni SMOTE) i
geometrijskih pro\v{s}irenja (G-SMOTE), \v{s}to ima i zna\v{c}ajnu
edukacijsku vrijednost.

\subsubsection{Arhitektura web su\v{c}elja}

Su\v{c}elje je realizirano kori\v{s}tenjem biblioteke Streamlit za
Python koja omogu\'{c}uje brzo prototipiranje interaktivnih web
aplikacija. Backend se izravno oslanja na implementirane SMOTE algoritme
iz paketa \texttt{smote\_variants}. Za vizualizaciju se koristi Plotly
(interaktivni scatter plotovi), dok se redukcija dimenzionalnosti provodi
algoritmima PCA i t-SNE iz biblioteke scikit-learn.

\subsubsection{Funkcionalnosti}

Opisuju se klju\v{c}ne funkcionalnosti su\v{c}elja:
u\v{c}itavanje vlastitog CSV skupa podataka ili odabir jednog od
ugra\dj{}enih skupova; redukcija dimenzionalnosti na 2D/3D pomo\'{c}u
PCA ili t-SNE; odabir jedne ili vi\v{s}e SMOTE varijanti za usporedbu;
pode\v{s}avanje parametara ($k$, postotak preuzorkovanja); usporedni
prikaz originalnih i sinteti\v{c}kih primjera; animacija procesa
generiranja za 2D podatke; prikaz granice odluke klasifikatora prije i
poslije preuzorkovanja te tablica metrika za kvantitativnu usporedbu.

\subsubsection{Prikazi web su\v{c}elja}

Prikazat \'{c}e se screenshotovi glavnih ekrana su\v{c}elja s kratkim
opisima: po\v{c}etni ekran s kontrolama, PCA prikaz originalnih podataka,
usporedni prikaz vi\v{s}e SMOTE varijanti, prikaz generiranja G-SMOTE u
odnosu na osnovni SMOTE, granica odluke prije i poslije preuzorkovanja
te tablica metrika.
